\subsection{A three-node ground query and the 663-bit combination}
\label{sec:e663}

The native stack equations permit replacing the four-node ground query
\(N(d,d)=d\) by the three-node query
\begin{equation}
                              Q(d)=d.
\label{eq:e663-query}
\end{equation}
This gives 688 bits when applied to the unchanged E692 equation list, or
\textbf{663 bits} when combined with the preceding 25-bit truth-equation
deletion.  The inspected input archive contains E692 and the four-node
query.  No independent priority claim is made for the elementary query
shortening.  The
exact 663-bit combination and the closure-based derivation of its omitted
axiom are the contribution being exhibited here.

To prove the equivalence directly, suppose the equation \(Q(d)=d\) is
added.  For arbitrary \(x,y\), the retained head and tail equations give
\begin{align*}
 x&=V(P(x,P(y,d)),d)\\
  &=V(P(x,P(y,d)),Q(d))\\
  &=V(P(y,d),d)=y.
\end{align*}
Thus the added ground equation entails collapse of the entire algebra.
Conversely, \(N(d,d)=d\) collapses the Boolean reduct furnished by Sh1:
it says \(\neg d=d\), whence
\(d=d\wedge d=d\wedge\neg d=0\) and
\(d=d\vee d=d\vee\neg d=1\).  Therefore every two elements coincide,
and in particular \(Q(d)=d\).

The same reasoning applies to the term algebra modulo equational
derivability, or equivalently uses completeness of equational logic.  For
the retained equation theory \(E\),
\[
 E\vdash Q(d)=d
 \quad\Longleftrightarrow\quad
 E\vdash N(d,d)=d.
\]
No intended meaning of the binder \(Q\), no raw-membership model, and no
truth value for the class \(d\) enters this argument.  The premises are
only the native head/tail equations and Sh1.  The chosen query is true
exactly in the trivial models of this theory.

The E663 accounting is
\[
               43+(7+4\cdot150+1)+4\cdot3=663.
\]
The new packet has 83 physical bytes with one zero padding bit.  Its
equations and ordered signature are exactly those of E667.  The script
emits E667 as a checkpoint and E663 as the shorter designated-query packet.
